\section{Target-Aware Observation Design}
\label{sec:design}

Once candidate values have been identified and estimated accurately enough to
distinguish them, the remaining task is to choose a feasible protocol. Because
the target variance $V_g(K)$ does not depend on the protocol, maximising the
normalised value $\cI_g(S;K)=F_g(S;K)/V_g(K)$ is equivalent to maximising
$F_g(S;K)$.
Exact rank-one marginal gains support forward selection with one-swap
refinement, and exhaustive search measures optimisation error when the
catalogue is small enough. Appendix~\ref{sec:app-design} gives the derivations,
algorithmic details and comparator formulas.

\subsection{Feasible Protocols and Design Objective}

For design, an acquisition action is represented by
\begin{equation}
a=(\ell_a,\;r_a,\;c_a),
\label{eq:action}
\end{equation}
where $\ell_a\in\R^p$ specifies the linear measurement, $r_a\ge0$ is its
effective noise variance, and $c_a>0$ is its cost. The observation is
\begin{equation}
Y_{i,a}=\ell_a^\top Z_i+\varepsilon_{i,a},
\qquad \varepsilon_{i,a}\sim\cN(0,r_a),
\label{eq:observation}
\end{equation}
with noises independent across actions and independent of $Z_i$. Timing and
support are encoded in $\ell_a$: a point measurement uses a coordinate vector,
whereas a window average uses normalised weights over its support. Repetition
is encoded through $r_a$; for example, averaging $M_a$ independent measurements
with per-measurement variance $\nu_a^2$ gives $r_a=\nu_a^2/M_a$. The cost may
include all repetitions and other acquisition burdens.

For a finite catalogue $\cV$ and cost budget $\mathcal B$, the feasible family
is
\begin{equation}
\Pi_{\mathcal B}=\Big\{S\subseteq\cV:\
\textstyle\sum_{a\in S}c_a\le\mathcal B\Big\}.
\label{eq:feasible-family}
\end{equation}
Application-specific restrictions, such as choosing at most one precision
variant on a given support, can be imposed on the catalogue or by taking a
subfamily of $\Pi_{\mathcal B}$. The results below apply to any such subfamily.
We restrict attention to feasible protocols for which every non-empty
$\Sigma_S(K)$ is nonsingular; this is automatic when every selected action has
$r_a>0$.

At the population level, observation design is therefore the optimisation problem
\begin{equation}
\max_{S\in\Pi_{\mathcal B}}\;F_g(S;K),
\label{eq:design-problem}
\end{equation}
where the candidate set $\cV$ ranges over acquisition times, window lengths,
repetition counts and noise levels. Given calibration data, the empirical
version replaces $K$ by $\widehat K$ and maximises $F_g(S;\widehat K)$;
\cref{cor:regret} bounds what that substitution costs.

\subsection{Exact Marginal Gains}

Suppressing the argument $K$ in $Q_S(K)$, write
\begin{equation}
P_S:=K-Q_S(K)=\Cov(Z\mid Y_S).
\label{eq:residual-covariance}
\end{equation}
Thus $v_{a\mid S}=P_S\ell_a$ is the covariance between a new measurement and
the residual trajectory, and $s_{a\mid S}$ is its conditional variance,
including measurement noise.

\begin{proposition}[Rank-one marginal gains]
\label{prop:marginal}
Let $S=\varnothing$, or assume that $\Sigma_S(K)$ is nonsingular. Let
$a\notin S$ satisfy
$S\cup\{a\}\in\Pi_{\mathcal B}$ and $s_{a\mid S}>0$, and put
\begin{equation}
v_{a\mid S}=P_S\ell_a,
\qquad
s_{a\mid S}=\ell_a^\top P_S\ell_a+r_a.
\label{eq:vs}
\end{equation}
Then $Q_{S\cup\{a\}}=Q_S+v_{a\mid S}v_{a\mid S}^\top/s_{a\mid S}$ and
$P_{S\cup\{a\}}=P_S-v_{a\mid S}v_{a\mid S}^\top/s_{a\mid S}$, so that the exact
marginal gain is
\begin{equation}
\Delta_g(a\mid S)
=\sum_{j,k}\omega_j\omega_k
\left[C_g\!\left\{Q_{S,jk}+\frac{v_{a\mid S,j}v_{a\mid S,k}}{s_{a\mid S}}\right\}
-C_g(Q_{S,jk})\right].
\label{eq:marginal-nonlinear}
\end{equation}
For the mean target this collapses to the closed form
\begin{equation}
\Delta_{\mathrm{mean}}(a\mid S)
=\frac{(\omega^\top P_S\ell_a)^2}{\ell_a^\top P_S\ell_a+r_a}.
\label{eq:marginal-mean}
\end{equation}
\end{proposition}

\Cref{eq:marginal-nonlinear} is the exact objective increment rather than a
surrogate. The rank-one form lets search update the residual covariance without
re-forming the full observation-covariance inverse.

The marginal-gain objective is monotone under nested protocols, although it
need not exhibit diminishing returns.
\begin{lemma}[Monotonicity under nested protocols]
\label{prop:monotone}
For every $g\in L^2(\phi)$, every correlation matrix $K$ with nonsingular
measurement covariance for each non-empty protocol concerned, and every
$S\subseteq S'$,
$F_g(S;K)\le F_g(S';K)$.
\end{lemma}

Adding measurements cannot decrease $F_g$. Monotonicity, however, does not
imply diminishing returns: even the mean target reduces to $R^2$ subset
selection, which is not submodular in general \citep{das2018submodular}.

For noiseless point measurements and a linear integral target, the objective
reduces to kernel quadrature
\citep{bach2017equivalence,briol2019probabilistic}. Measurement noise and
nonlinear targets place the general problem outside that noiseless formulation.

\subsection{Search and Comparators}
\label{sec:algorithms}

Starting from $S=\varnothing$ and $P_S=K$, forward selection repeatedly adds
the feasible action with largest exact gain in
\eqref{eq:marginal-nonlinear}, using gain per unit cost when costs differ, and
then applies the rank-one update. A one-swap refinement makes feasible
selected--unselected exchanges whenever the exchange improves the objective.
Because monotonicity does not imply submodularity, exhaustive search is used as an
optimisation benchmark when the catalogue is small enough. Pseudocode and
operation counts are given in Appendix~\ref{sec:app-design-search}.

To isolate which information drives a schedule, we compare the target-aware
objective with several alternatives. Latent-state
mutual information maximises $I(Z;Y_S)$ and is target-free. Integrated
posterior variance minimises $\sum_j\omega_j(P_S)_{jj}$, so it uses temporal
weights but not the target transformation. Linear-target design retains the
actual weights and noise but sets $g(z)=z$, thereby maximising
$\omega^\top Q_S(K)\omega$. Noiseless kernel quadrature uses the same linear
criterion with $R_S=0$. These comparisons separate target nonlinearity from
the noise model and from target-free coverage criteria; their exact objectives
and marginal gains are recorded in Appendix~\ref{sec:app-design-comparators}.
